%% kapa: exam=1 round=30 year=2019 period=2 form=A subject=accounting subject_ko=회계학 seq=22 number=62
%% source: Q-Net 감정평가사 기출문제 (cst00309) artlSeq=5205205 file=extracted/1008/5205205/2019년 제30회 감정평가사 1차시험 문제지/2019년 제30회 감정평가사 1차시험 문제지/2019년도 제30회 감정평가사 1차 2교시 A형.hwp
%% conversion: kordoc 4.24.51 -> tools/gen1.py (auto); review=proofread-sonnet; flags=none
\begin{question}{62}
금융상품에 관한 설명으로 옳은 것은?
\begin{choices}
  \item 당기손익-공정가치로 측정되는 ‘지분상품에 대한 특정 투자’에 대해서는 후속적인 공정가치 변동은 최초 인식시점이라 하더라도 기타포괄손익으로 표시하도록 선택할 수 없다.
  \item 측정이나 인식의 불일치, 즉 회계불일치의 상황이 아닌 경우 금융자산은 금융자산의 관리를 위한 사업모형과 금융자산의 계약상 현금흐름의 특성 모두에 근거하여 상각후원가, 기타포괄손익-공정가치, 당기손익-공정가치로 측정되도록 분류한다.
  \item 금융자산 전체나 일부의 회수를 합리적으로 예상할 수 없는 경우에도 해당 금융자산의 총 장부금액을 직접 줄일 수는 없다.
  \item 기타포괄손익-공정가치 측정 금융자산의 기대신용손실을 조정하기 위한 기대신용손실액(손상차손)은 당기손실로 인식하고, 기대신용손실환입액(손상차손환입)은 기타포괄손익으로 인식한다.
  \item 금융자산을 상각후원가 측정범주에서 기타포괄손익-공정가치 측정 범주로 재분류하는 경우 재분류일의 공정가치로 측정하며, 재분류 전 상각후원가와 공정가치 차이에 따른 손익은 당기손익으로 인식한다.
\end{choices}
\end{question}
\answer{2}
